\section{二分类的概率建模与对数损失}\label{ch:logistic}

\subsection{从线性得分到条件概率}

考虑二分类标签 $Y\in\{0,1\}$，线性得分 $z=x\trans\beta$ 可取任意实数，因此不能直接作为概率。逻辑回归使用 Logistic 函数
\begin{equation}\label{eq:sigmoid}
\sigma(z)=\frac{1}{1+\mathrm e^{-z}}
\end{equation}
把实数映射到开区间 $(0,1)$，并定义
\begin{equation}\label{eq:logistic-model}
\pi(x)=\Prob(Y=1\mid X=x)=\sigma(x\trans\beta).
\end{equation}
其中 $x$ 可以含截距对应的常数分量。这里采用 $\pi$ 表示条件概率，以免与特征维数 $p$ 混淆。

由 \eqref{eq:sigmoid} 可得
\[
\frac{\pi(x)}{1-\pi(x)}=\mathrm e^{x\trans\beta},
\]
因此
\begin{equation}\label{eq:log-odds}
\boxed{\log\frac{\pi(x)}{1-\pi(x)}=x\trans\beta.}
\end{equation}
模型对\emph{对数几率}作线性建模。这也解释了 Logistic Regression 虽然沿用“回归”一词，却主要用于分类任务\cite[第3章]{zhou2016}。

\subsection{Sigmoid 的导数}

对式 \eqref{eq:sigmoid} 求导，
\begin{align}
\sigma'(z)
&=-(1+\mathrm e^{-z})^{-2}(-\mathrm e^{-z})\\
&=\frac{\mathrm e^{-z}}{(1+\mathrm e^{-z})^2}\\
&=\boxed{\sigma(z)[1-\sigma(z)]}.\label{eq:sigmoid-derivative}
\end{align}
因此 Sigmoid 单调递增，且 $0<\sigma'(z)\le1/4$。\textbf{Sigmoid 并不是在整个实数轴上凸的}：其二阶导数符号在 $z=0$ 两侧发生改变。逻辑回归的凸性来自负对数似然关于参数的结构，而非 Sigmoid 函数自身的全局凸性。

\begin{figure}[htbp]
\centering
\begin{tikzpicture}
\begin{axis}[width=.74\textwidth,height=6cm,xlabel={$z$},ylabel={$\sigma(z)$},xmin=-6,xmax=6,ymin=0,ymax=1.05,grid=major,grid style={draw=gray!16},ytick={0,0.25,0.5,0.75,1},axis line style={draw=black!55}]
\addplot[very thick,color=accent,domain=-6:6,samples=160]{1/(1+exp(-x))};
\addplot[densely dashed,color=black!45] coordinates {(-6,.5) (6,.5)};
\end{axis}
\end{tikzpicture}
\caption{Logistic 函数将线性得分映射为二分类概率。$z=0$ 时输出为 $0.5$。}\label{fig:sigmoid}
\end{figure}

\subsection{伯努利似然与交叉熵}

对样本 $(x_i,y_i)$，设 $Y_i\mid x_i\sim\operatorname{Bernoulli}(\pi_i)$，其中 $\pi_i=\sigma(x_i\trans\beta)$。在给定输入后，若样本标签条件独立，则条件似然为
\begin{equation}\label{eq:logit-likelihood}
\mathcal L(\beta)=\prod_{i=1}^{m}\pi_i^{y_i}(1-\pi_i)^{1-y_i}.
\end{equation}
取对数、乘以负号并按样本平均，得到
\begin{equation}\label{eq:bce}
\boxed{J(\beta)=-\frac1m\sum_{i=1}^{m}\left[y_i\log\pi_i+(1-y_i)\log(1-\pi_i)\right].}
\end{equation}
这是二元交叉熵，也是该伯努利模型的平均负对数似然。由此可见，在给定概率假设和独立性条件下，损失函数具有明确的统计来源，而不是仅凭曲线形状选择\cite[第6章]{lihang2019}。

\subsection{从链式法则推导梯度}

先考虑一个样本，记
\[
\ell=-y\log\pi-(1-y)\log(1-\pi),\quad\pi=\sigma(z),\quad z=x\trans\beta.
\]
对概率变量求导：
\[
\frac{\partial\ell}{\partial\pi}=-\frac y\pi+\frac{1-y}{1-\pi}
=\frac{\pi-y}{\pi(1-\pi)}.
\]
再利用式 \eqref{eq:sigmoid-derivative}，
\[
\frac{\partial\ell}{\partial z}
=\frac{\partial\ell}{\partial\pi}\frac{\mathrm d\pi}{\mathrm dz}
=\frac{\pi-y}{\pi(1-\pi)}\pi(1-\pi)=\pi-y.
\]
由于 $\nabla_\beta z=x$，因此
\begin{equation}\label{eq:logistic-single-gradient}
\boxed{\nabla_\beta\ell=(\pi-y)x.}
\end{equation}
把所有样本堆叠为设计矩阵 $X\in\R^{m\times p}$，定义概率向量 $\boldsymbol\pi=(\pi_1,\ldots,\pi_m)\trans$，得到平均损失梯度
\begin{equation}\label{eq:logistic-full-gradient}
\boxed{\nabla J(\beta)=\frac1mX\trans(\boldsymbol\pi-y).}
\end{equation}
这里的结构与最小二乘梯度 \eqref{eq:ls-gradient} 十分接近：都是先计算某种“预测与真实值之差”，再乘设计矩阵的转置。差别在于预测值和损失函数来自不同的概率或建模假设。

\subsection{Hessian、凸性与例外情形}

记 $W=\operatorname{diag}\{\pi_i(1-\pi_i)\}_{i=1}^m$，则
\begin{equation}\label{eq:logistic-hessian}
\nabla^2J(\beta)=\frac1mX\trans W X.
\end{equation}
对任意 $v$，
\[
v\trans\nabla^2J(\beta)v=\frac1m(Xv)\trans W(Xv)\ge0,
\]
所以在常规无正则化二分类逻辑回归中，平均负对数似然关于 $\beta$ 是凸函数。对所有有限参数，$\pi_i\in(0,1)$，故 $W$ 的对角元素严格为正；若 $X$ 满列秩，则 Hessian 正定。

但是正定 Hessian 并不自动保证有限的极小点存在。例如样本线性完全可分时，沿某些参数方向把系数绝对值持续放大，可以使负对数似然逼近下确界 $0$，而无正则化极大似然估计可能不存在有限解。因而“凸”与“存在唯一有限最优解”是不同的判断。

\subsection{一次相容的手工更新}

设单个样本 $x=1$、$y=1$，初始 $w_0=b_0=0$，模型为 $\pi=\sigma(wx+b)$。此时 $z=0$、$\pi=1/2$，因此
\[
\frac{\partial\ell}{\partial w}=(1/2-1)\cdot 1=-1/2,\qquad
\frac{\partial\ell}{\partial b}=-1/2.
\]
取学习率 $\eta=0.1$，同步更新后
\begin{equation}
(w_1,b_1)=(0.05,0.05).
\end{equation}
这个例子特意从给定参数重新计算概率，避免出现“人为指定的概率值”与 Sigmoid 模型不一致的问题。概率预测若需要转为类别，还须另外声明分类阈值及相应的决策代价。

\subsection{模型解释与适用范围}

在其他特征固定时，某个系数 $\beta_j$ 增加一单位特征对应的对数几率变化为 $\beta_j$；对应的几率乘以 $\mathrm e^{\beta_j}$。这一解释针对模型规定的条件几率，并不能仅凭相关性就作因果推断。样本不平衡、特征高度相关、概率校准等问题留待后续专题讨论。
